\section{Introduction}
\label{sec:intro}

\IEEEPARstart{A}{} brain--computer interface (BCI) reads activity from the brain and turns it
into an instruction for a machine, without passing through the nerves and muscles that
normally carry a movement~\cite{wolpaw2002}. The most common non-invasive form records the
electroencephalogram (EEG) from the scalp. In a motor-imagery (MI) BCI the user only
imagines a movement. Imagining a movement weakens the mu and beta rhythms over the motor
cortex, an effect called event-related desynchronisation
(ERD)~\cite{pfurtscheller1999,pfurtscheller2001}, and a decoder is trained to recognise
that change. Two lines of research use such interfaces. In stroke rehabilitation the BCI
closes the loop between a patient's intention and a movement made by a device, and
controlled trials suggest that this aids recovery of the arm~\cite{cervera2018,biasiucci2018}.
In motor augmentation, healthy people learn to operate an extra robotic digit worn on the
hand~\cite{kieliba2021,eden2022}.

Wang et al.\ brought the two lines together. They asked healthy
volunteers~\cite{wang2024jbhi} and later stroke patients~\cite{wang2024tnsre} to imagine
flexing a robotic sixth finger strapped to the wrist, and they released a dataset of
\nSubjects{} patients who imagined either this sixth finger or a grasp of the affected hand
over three sessions~\cite{wang2026dataset}. Because the dataset is public, the paradigm can
now be studied by groups who did not record it, and that is what we do here.

Every MI-BCI needs a calibration step before it can be used: the user performs a series of
labelled trials from which the decoder learns that person's pattern~\cite{lotte2015}. For
a stroke patient, each minute spent on calibration is effort taken away from therapy, and
the step is usually repeated every day because the EEG drifts. Transfer learning reuses
recordings of other people and is known to shorten calibration for healthy users who
imagine familiar movements~\cite{zanini2018,he2020,wu2022review}. Whether it helps when the
users are stroke patients and the command is one that no person has ever performed, a sixth
finger, has not been tested.

A study of calibration first needs an evaluation that does not leak. The dataset paper
reports accuracies of \rqOnePublishedGroupBCspLda\,\% and \rqOnePublishedGroupTCspLda\,\%
for telling imagined hand from imagined sixth-finger movements before therapy, in its two
therapy groups~\cite{wang2026dataset}. Its released analysis script fits the common spatial
pattern (CSP) filters on all the trials of a session and only afterwards splits the trials
for cross-validation. Letting test trials shape a supervised step in this way is a known
cause of optimistic estimates~\cite{kapoor2023,brookshire2024}. We therefore re-evaluated
the dataset with every fitted step kept inside the training data, and every later result
in this paper rests on that evaluation.

The study asks four research questions (RQs). RQ1 measures how much the published protocol
inflates accuracy on these recordings, and whether the same code inflates accuracy on two
healthy datasets. RQ2 asks how much calibration a new patient needs before a decoder can
tell rest, the affected hand and the sixth finger apart, and whether other patients' data
reduce it. RQ3 asks the same when the patient's hand trials already exist and only the
sixth finger has to be added. RQ4 gathers exploratory questions: continuous use of the
decoder, the side of the lesion, and use on later days. The hypotheses for RQ1 to RQ3 were
written down and frozen before the main study. The contributions are:
\begin{enumerate}
  \item a benchmark of this sixth-finger stroke dataset that keeps every fitted step away from
        the test data, preceded by an audit that uncovered duplicated recordings;
  \item what we believe is the first study of how much calibration a supernumerary, novel
        motor command needs in stroke patients, with hypotheses fixed in advance for two
        cases: a new user, and a new command added once the hand trials exist;
  \item Riemannian and deep transfer compared in this population, with analyses specific to
        stroke (mirroring the lesioned hemisphere, later sessions) kept exploratory;
  \item an ablation study that collects every component comparison the project had run,
        each tied to its results file and labelled as pre-registered or exploratory; and
  \item a demonstrator that streams held-out recordings through a decoder, on a server or in
        the browser, and animates a virtual hand wearing a sixth finger.
\end{enumerate}

Sections~\ref{sec:lit} to~\ref{sec:conclusion} cover, in turn, the earlier work, the
methodology, the architecture, the pre-registered experiments, the ablation study, the
discussion with its limitations and future scope, and the conclusion.
